%% kapa: exam=1 round=22 year=2011 period=2 form=A subject=appraisal_law subject_ko=부동산관계법규 seq=31 number=71
%% source: Q-Net 감정평가사 기출문제 (cst00309) artlSeq=5036790 file=extracted/1008/5036790/2011년 제22회 감정평가사 1차시험 기출문제/2교시.hwp
%% conversion: kordoc 4.24.51 -> tools/gen1.py (auto); review=proofread-gpt-astra; flags=none
\begin{question}{71}
건축법령상 건축물의 높이제한에 관한 설명으로 옳지 않은 것은?
\begin{choices}
  \item 건축물의 높이제한은 주로 건축물의 안전 확보, 일조, 통풍, 채광 등 위생적이고 쾌적한 주거환경의 확보 및 도시미관 유지 등을 도모하기 위한 것이다.
  \item 건축법상 일조권의 확보를 위한 건축물의 높이를 제한하는 지역은 원칙적으로 전용주거지역, 일반주거지역 및 준주거지역이다.
  \item 최고 높이가 정하여지지 않은 가로구역의 경우 건축물의 각 부분의 높이는 원칙적으로 그 부분으로부터 전면도로의 반대쪽 경계선까지의 수평거리의 1.5배를 넘을 수 없다.
  \item 허가권자는 같은 가로구역에서 건축물의 용도 및 형태에 따라 건축물의 높이를 다르게 정할 수 있다.
  \item 특별자치도지사 또는 시장ㆍ군수ㆍ구청장은 필요하다고 인정하는 경우 가로구역의 최고 높이를 대통령령이 정하는 바에 따라 건축위원회의 심의를 거쳐 완화하여 적용할 수 있다.
\end{choices}
\end{question}
\answer{2}
